% !TEX root = ../discretemath.tex

\section{Дискретная теория вероятности}

\subsection{Конечное вероятностное пространство}

\begin{Definition}{Конечное вероятностное пространство}{}
	Пусть \(U\) --- конечное множество элементарных исходов.
	
	Функция
	\[
	P_r \colon U \to [0,1]
	\]
	называется \textit{вероятностным распределением} на \(U\), если
	\[
	\sum_{\omega \in U} P_r(\omega)=1.
	\]
	
	Тогда пара
	\[
	(U, P_r)
	\]
	называется \textit{конечным вероятностным пространством}.
\end{Definition}

\begin{Definition}{Событие и его вероятность}{}
	Пусть \((U, P_r)\) --- конечное вероятностное пространство.
	
	Любое подмножество \(A \subseteq U\) называется \textit{событием}.
	
	Вероятность события \(A\) определяется формулой
	\[
	P_r(A)=\sum_{\omega\in A} P_r(\omega).
	\]
\end{Definition}

\begin{Statement}{Свойства вероятности}{}
	Пусть \(A,B \subseteq U\). Тогда:
	\begin{enumerate}[label=\arabic*.]
		\item \(P_r(\varnothing)=0,\qquad P_r(U)=1.\)
		\item \(P_r(A)+P_r(U\setminus A)=1.\)
		\item \(P_r(A\cup B)=P_r(A)+P_r(B)-P_r(A\cap B).\)
	\end{enumerate}
\end{Statement}
	\subsection{Примеры}
	
	\begin{Example}{$n$ подбрасываний монеты}{}
		Рассмотрим \(n\) подбрасываний честной монеты. Тогда каждый исход --- это последовательность из \(n\) символов, где, например, \(1\) означает орла, а \(0\) --- решку.
		
		Найдём вероятность того, что выпало ровно \(k\) орлов.
		\end{Example}
		Множество исходов:
		\[
		U=\{0,1\}^n,
		\qquad |U|=2^n.
		\]
		
		Так как монета честная, каждый исход равновероятен:
		\[
		P_r(\omega)=2^{-n}, \qquad \omega\in U.
		\]
		
	   Пусть
		\[
		A_k=\{\omega\in U \mid |\omega|=k\},
		\]
		где \(|\omega|\) --- число орлов в последовательности \(\omega\).
		
		Тогда
		\[
		P_r(A_k)=\sum_{\omega\in A_k}P_r(\omega)
		=\sum_{\substack{\omega\in U\\ |\omega|=k}}2^{-n}.
		\]
		Так как таких последовательностей ровно \(\binom{n}{k}\), получаем
		\[
		P_r(\text{ровно }k\text{ орлов})
		=
		\binom{n}{k}2^{-n}.
		\]

	
	\begin{Example}{7 карт: вероятность того, что есть 5 карт одной масти}{}
		Из стандартной колоды в \(52\) карты случайно выбираются \(7\) карт. Требуется найти вероятность того, что среди них есть \textit{ровно} 5 карт одной масти, а затем учесть выбор масти.
	\end{Example}
		
		Сначала выберем масть: это можно сделать \(4\) способами.
		
		Далее:
		\begin{itemize}
			\item 5 карт этой масти выбираются из 13:
			\[
			\binom{13}{5};
			\]
			\item оставшиеся 2 карты должны быть не этой масти, то есть выбираются из \(39\) карт:
			\[
			\binom{39}{2}.
			\]
		\end{itemize}
		
		Общее число наборов из 7 карт:
		\[
		\binom{52}{7}.
		\]
		
		Следовательно,
		\[
		P(\text{ровно 5 карт одной масти})
		=
		\frac{4\binom{13}{5}\binom{39}{2}}{\binom{52}{7}}.
		\]


	
	\begin{Example}{Задача о потерянном билете}{}
		В автобус заходит \(N\) человек. У первого пассажира билет потерян.
		
		\begin{itemize}
			\item Первый пассажир садится на \textit{случайное} место.
			\item Каждый следующий пассажир:
			\begin{itemize}
				\item садится на своё место, если оно свободно;
				\item иначе садится на случайное из оставшихся свободных мест.
			\end{itemize}
		\end{itemize}
		
		Найти вероятность того, что \(N\)-й пассажир сядет на своё место.
	\end{Example}
	
	\subsubsection*{Решение}
	
	Обозначим через \(a_n\) вероятность того, что в описанной модели \(n\)-й пассажир в итоге сядет на своё место.
	
	Нужно найти \(a_n\).
	
	После первого шага первый пассажир случайно выбирает одно из \(n\) мест.
	
	Разберём случаи:
	
	\begin{itemize}
		\item Если он сразу сел на место \(1\), то все остальные сядут на свои места, значит вклад в вероятность равен
		\[
		\frac1n\cdot 1.
		\]
		
		\item Если он сел на место \(n\), то \(n\)-й пассажир уже не сможет сесть на своё место, вклад равен
		\[
		\frac1n\cdot 0.
		\]
		
		\item Если он сел на место \(k\), где \(2\le k\le n-1\), то пассажиры \(2,3,\dots,k-1\) сядут на свои места, а пассажир \(k\) окажется в точно такой же задаче, но уже для \(n-k+1\) пассажиров:
		\[
		\frac1n\cdot a_{n-k+1}.
		\]
	\end{itemize}
	
	Следовательно,
	\[
	a_n
	=
	\frac1n\left(1+a_{n-1}+a_{n-2}+\dots+a_2+0\right).
	\]
	
	То есть
	\[
	a_n=\frac1n\left(1+\sum_{j=2}^{n-1}a_j\right).
	\]
	
	\subsection*{База}
	
	При \(n=2\):
	\[
	a_2=\frac12,
	\]
	потому что первый пассажир с равной вероятностью садится либо на место \(1\), либо на место \(2\).
	
	При \(n=3\) аналогично получаем
	\[
	a_3=\frac12.
	\]
	
	\begin{Proposition}{}{}
		Для всех \(n\ge 2\) выполнено
		\[
		a_n=\frac12.
		\]
	\end{Proposition}
	
	\begin{proof}
		Докажем по индукции.
		
		\textbf{База:} для \(n=2\)
		\[
		a_2=\frac12.
		\]
		
		\textbf{Переход:} предположим, что
		\[
		a_2=a_3=\dots=a_{n-1}=\frac12.
		\]
		Тогда из рекуррентной формулы получаем
		\[
		a_n
		=
		\frac1n\left(1+\sum_{j=2}^{n-1}a_j\right)
		=
		\frac1n\left(1+\sum_{j=2}^{n-1}\frac12\right).
		\]
		Слагаемых в сумме \((n-2)\), значит
		\[
		a_n
		=
		\frac1n\left(1+\frac{n-2}{2}\right)
		=
		\frac1n\cdot\frac{n}{2}
		=
		\frac12.
		\]
		
		Следовательно,
		\[
		a_n=\frac12
		\]
		для всех \(n\ge 2\).
	\end{proof}


\begin{Example}{}{}
	Пусть в группе 40 человек: какова вероятность, что у двух совпадает день рождения?
\end{Example}

Удобнее сначала рассмотреть противоположное событие: \textit{у всех дни рождения различны}.

Общее число способов назначить дни рождения \(40\) людям:
\[
365^{40}.
\]

Число благоприятных исходов для события ``все дни рождения различны'' равно
\[
365\cdot 364\cdot 363\cdots 326
=
\frac{365!}{325!} < (345.5)^{40}.
\]

поскольку среднее первого и последнего множителя равно
\[
\frac{365+326}{2}=345.5,
\]
а произведение не превосходит степени среднего арифметического:
\[
ab\le \left(\frac{a+b}{2}\right)^2.
\]

Следовательно,
\[
P(\text{у всех разные дни рождения})
=
\frac{365\cdot 364\cdot \ldots \cdot 326}{365^{40}} <
\left(\frac{345.5}{365}\right)^{40}
=
\left(1-\frac{19.5}{365}\right)^{40}.
\]

Значит,
\[
P(\text{у двух есть совпадающий день рождения})
=
1-\frac{365\cdot 364\cdot \ldots \cdot 326}{365^{40}}.
\]


\begin{Theorem}{Формула включений--исключений}{}
	Пусть \(A_1,\dots,A_k\subseteq U\). Тогда
	\[
	P_r\!\left(\bigcup_{i=1}^k A_i\right)
	=
	\sum_{\varnothing\neq I\subseteq \{1,\dots,k\}}
	(-1)^{|I|-1}
	P_r\!\left(\bigcap_{i\in I} A_i\right).
	\]
\end{Theorem}

То есть в развёрнутом виде:
\[
P_r(A_1\cup\cdots\cup A_k)
=
\sum_{i=1}^k P_r(A_i)
-
\sum_{1\le i<j\le k}P_r(A_i\cap A_j)
+\cdots
+(-1)^{k-1}P_r(A_1\cap\cdots\cap A_k).
\]

\begin{Example}{Случайная перестановка без неподвижных точек}{}
	Рассмотрим случайную перестановку \(\sigma\in S_n\), равновероятно выбранную из всех \(n!\) перестановок.
	
	Какова вероятность того, что у \(\sigma\) нет неподвижных точек?
	То есть что
	\[
	\sigma(i)\ne i
	\qquad \text{для всех } i=1,\dots,n.
	\]
\end{Example}

Для каждого \(i=1,\dots,n\) введём событие
\[
B_i=\{\sigma\in S_n \mid \sigma(i)=i\}.
\]
Тогда событие ``нет неподвижных точек'' есть
\[
\overline{B_1}\cap \cdots \cap \overline{B_n}.
\]

Следовательно,
\[
P_r(\text{нет неподвижных точек})
=
1-P_r(B_1\cup\cdots\cup B_n).
\]

По формуле включений--исключений:
\[
P_r(\text{нет неподвижных точек})
=
1+
\sum_{\varnothing\neq I\subseteq \{1,\dots,n\}}
(-1)^{|I|}
P_r\!\left(\bigcap_{i\in I} B_i\right).
\]

Теперь вычислим
\[
P_r\!\left(\bigcap_{i\in I} B_i\right).
\]
Если фиксированы все точки из множества \(I\), то остальные \(n-|I|\) элементов можно переставить произвольно. Поэтому число таких перестановок равно
\[
(n-|I|)!,
\]
а значит,
\[
P_r\!\left(\bigcap_{i\in I} B_i\right)
=
\frac{(n-|I|)!}{n!}.
\]

Подставляя, получаем
\[
P_r(\text{нет неподвижных точек})
=
1+
\sum_{\varnothing\neq I\subseteq \{1,\dots,n\}}
(-1)^{|I|}
\frac{(n-|I|)!}{n!}.
\]

Сгруппируем слагаемые по \(k=|I|\). Подмножеств мощности \(k\) ровно \(\binom{n}{k}\), поэтому
\[
P_r(\text{нет неподвижных точек})
=
1+\sum_{k=1}^n \binom{n}{k}(-1)^k\frac{(n-k)!}{n!}.
\]

Но
\[
\binom{n}{k}\frac{(n-k)!}{n!}
=
\frac{n!}{k!(n-k)!}\cdot \frac{(n-k)!}{n!}
=
\frac1{k!}.
\]
Значит,
\[
P_r(\text{нет неподвижных точек})
=
1+\sum_{k=1}^n (-1)^k\frac1{k!}.
\]

Итак,
\[
P_r(\text{нет неподвижных точек})
=
1-\frac1{1!}+\frac1{2!}-\frac1{3!}+\cdots+(-1)^n\frac1{n!}.
\]

\begin{Remark}{}{}
При \(n\to\infty\)
\[
P_r(\sigma \text{ не имеет неподвижных точек})
\longrightarrow
\frac1e.
\]
\end{Remark}

\subsection{Условная вероятность}

\begin{Definition}{Условная вероятность}{}
	Пусть \(A,B\subseteq U\) и \(P_r(B)>0\). Тогда \textit{условной вероятностью события \(A\) при условии \(B\)} называется величина
	\[
	P_r(A\mid B)=\frac{P_r(A\cap B)}{P_r(B)}.
	\]
\end{Definition}

\subsection*{Свойства условной вероятности}

\begin{Statement}{}{}
	Пусть \(P_r(C)>0\). Тогда для любых событий \(A,B\) выполняется
	\[
	P_r(A\cup B\mid C)
	=
	P_r(A\mid C)+P_r(B\mid C)-P_r(A\cap B\mid C).
	\]
\end{Statement}

\begin{proof}
	По определению условной вероятности
	\[
	P_r(A\cup B\mid C)
	=
	\frac{P_r((A\cup B)\cap C)}{P_r(C)}.
	\]
	Но
	\[
	(A\cup B)\cap C=(A\cap C)\cup(B\cap C),
	\]
	поэтому
	\[
	P_r((A\cup B)\cap C)
	=
	P_r(A\cap C)+P_r(B\cap C)-P_r(A\cap B\cap C).
	\]
	Делим на \(P_r(C)\):
	\[
	P_r(A\cup B\mid C)
	=
	\frac{P_r(A\cap C)}{P_r(C)}
	+
	\frac{P_r(B\cap C)}{P_r(C)}
	-
	\frac{P_r(A\cap B\cap C)}{P_r(C)}.
	\]
	Отсюда и следует формула:
	\[
	P_r(A\cup B\mid C)
	=
	P_r(A\mid C)+P_r(B\mid C)-P_r(A\cap B\mid C).
	\]
\end{proof}

\begin{Statement}{}{}
	Пусть \(P_r(B\cap C)>0\) и \(P_r(C)>0\). Тогда
	\[
	P_r(A\mid B\cap C)=\frac{P_r(A\cap B\mid C)}{P_r(B\mid C)}.
	\]
\end{Statement}

\begin{proof}
	По определению
	\[
	P_r(A\mid B\cap C)=\frac{P_r(A\cap B\cap C)}{P_r(B\cap C)}.
	\]
	С другой стороны,
	\[
	P_r(A\cap B\mid C)=\frac{P_r(A\cap B\cap C)}{P_r(C)},
	\qquad
	P_r(B\mid C)=\frac{P_r(B\cap C)}{P_r(C)}.
	\]
	Поэтому
	\[
	\frac{P_r(A\cap B\mid C)}{P_r(B\mid C)}
	=
	\frac{\dfrac{P_r(A\cap B\cap C)}{P_r(C)}}{\dfrac{P_r(B\cap C)}{P_r(C)}}
	=
	\frac{P_r(A\cap B\cap C)}{P_r(B\cap C)}
	=
	P_r(A\mid B\cap C).
	\]
\end{proof}

\begin{Example}{Две монеты}{}
	Рассмотрим два подбрасывания честной монеты.
	
	Пусть
	\[
	A=\{\text{выпало два орла}\},
	\qquad
	B=\{\text{на первой монете выпал орёл}\}.
	\]
	Тогда
	\[
	P_r(A)=\frac14,
	\qquad
	P_r(B)=\frac12.
	\]
	Поскольку событие \(A\) влечёт \(B\), имеем
	\[
	A\cap B=A,
	\]
	и значит
	\[
	P_r(A\mid B)=\frac{P_r(A\cap B)}{P_r(B)}
	=
	\frac{\frac14}{\frac12}
	=
	\frac12.
	\]
\end{Example}

\subsection{Независимость событий}

\begin{Definition}{Независимость двух событий}{}
	События \(A\) и \(B\) называются \textit{независимыми}, если
	\[
	P_r(A\cap B)=P_r(A)\,P_r(B).
	\]
\end{Definition}

\begin{Remark}{}{}
	Если \(P_r(B)=0\), то для любого события \(A\)
	\[
	P_r(A\cap B)=0=P_r(A)\cdot P_r(B),
	\]
	поэтому \(A\) и \(B\) автоматически независимы.
\end{Remark}

\begin{Example}{Кубик}{}
	Бросается честный кубик, множество исходов:
	\[
	U=\{1,2,3,4,5,6\}.
	\]
	
	Пусть
	\[
	A=\{\text{выпало число } \ge 4\}=\{4,5,6\},
	\]
	\[
	B=\{\text{выпало число, делящееся на }3\}=\{3,6\}.
	\]
\end{Example}

	Тогда
	\[
	P_r(A)=\frac36=\frac12,
	\qquad
	P_r(B)=\frac26=\frac13.
	\]
	Пересечение:
	\[
	A\cap B=\{6\},
	\]
	поэтому
	\[
	P_r(A\cap B)=\frac16.
	\]
	С другой стороны,
	\[
	P_r(A)\,P_r(B)=\frac12\cdot\frac13=\frac16.
	\]
	Следовательно, события \(A\) и \(B\) независимы.


\subsection{Независимость от набора событий}

\begin{Definition}{Независимость события от семейства событий}{}
	Пусть даны событие \(A\) и события \(B_1,\dots,B_n\).
	
	Говорят, что событие \(A\) \textit{независимо от событий \(B_1,\dots,B_n\)}, если для любого множества индексов $I\subseteq \{1,\dots,n\}$ события $A \quad \text{и} \quad \bigcap_{i\in I} B_i$ независимы.
\end{Definition}

\begin{Definition}{Независимость в совокупности}{}
	События \(A_1,\dots,A_n\) называются \textit{независимыми в совокупности}, если для любого множества индексов
	\[
	I\subseteq \{1,\dots,n\}
	\]
	выполняется
	\[
	P_r\!\left(\bigcap_{i\in I} A_i\right)
	=
	\prod_{i\in I} P_r(A_i).
	\]
\end{Definition}

\begin{Example}{Попарная независимость без независимости в совокупности}{}
	Рассмотрим \(n\) независимых подбрасываний честной монеты.
	
	Для \(i=1,\dots,n\) обозначим через \(B_i\) событие
	\[
	B_i=\{\text{на \(i\)-м подбрасывании выпал орёл}\}.
	\]
	Кроме того, рассмотрим событие
	\[
	B_{n+1}=\{\text{общее число орлов чётно}\}.
	\]
	
	Непосредственно проверяется, что любые два события из набора
	\[
	B_1,\dots,B_n,B_{n+1}
	\]
	независимы.
	
	Однако события
	\[
	B_1,\dots,B_n,B_{n+1}
	\]
	не являются независимыми в совокупности, поскольку событие \(B_{n+1}\)
	однозначно определяется по событиям \(B_1,\dots,B_n\).
\end{Example}

\begin{Statement}{}{}
	Любые \(n\) событий из набора
	\[
	B_1,\dots,B_n,B_{n+1}
	\]
	независимы в совокупности.
\end{Statement}

\begin{proof}
	Нужно рассмотреть два случая.
	
	\textbf{1) Выбраны только события из \(B_1,\dots,B_n\).}
	
	Тогда это просто независимость результатов отдельных подбрасываний монеты.
	
	\textbf{2) Среди выбранных событий есть \(B_{n+1}\), но какого-то одного из событий \(B_1,\dots,B_n\) нет.}
	
	НУО отсутствует событие \(B_n\). Тогда рассматривается семейство
	\[
	B_1,\dots,B_{n-1},B_{n+1}.
	\]
	
	Зафиксируем любые значения первых \(n-1\) монет. После этого исход \(n\)-й монеты остаётся случайным и равновероятным. Среди двух возможных исходов \(n\)-й монеты ровно один делает общее число орлов чётным. Поэтому
	\[
	P_r(B_{n+1}\mid \text{фиксированы }B_1,\dots,B_{n-1})=\frac12.
	\]
	Но и
	\[
	P_r(B_{n+1})=\frac12.
	\]
	Значит, событие \(B_{n+1}\) независимо от любой комбинации событий \(B_1,\dots,B_{n-1}\). Следовательно,
	\[
	B_1,\dots,B_{n-1},B_{n+1}
	\]
	независимы в совокупности.
	
	Точно так же рассматривается любой другой набор из \(n\) событий.
\end{proof}

\subsection{Формула Байеса}

\begin{Statement}{Формула Байеса}{}
	Пусть \(P_r(A)>0\) и \(P_r(B)>0\). Тогда
	\[
	P_r(A\mid B)=\frac{P_r(B\mid A)\,P_r(A)}{P_r(B)}.
	\]
\end{Statement}

\begin{proof}
	По определению условной вероятности
	\[
	P_r(A\mid B)=\frac{P_r(A\cap B)}{P_r(B)},
	\qquad
	P_r(B\mid A)=\frac{P_r(A\cap B)}{P_r(A)}.
	\]
	Из второй формулы получаем
	\[
	P_r(A\cap B)=P_r(B\mid A)\,P_r(A).
	\]
	Подставляя это в первую, имеем
	\[
	P_r(A\mid B)=\frac{P_r(B\mid A)\,P_r(A)}{P_r(B)}.
	\]
\end{proof}